\chapter{Reproducción del caso de uso de validación}\label{ap:reproduccion}

Este anexo concentra los comandos exactos para reproducir el caso de uso de validación desde un repositorio limpio, clonado del repositorio público del TFM en GitHub.\footnote{\url{https://github.com/AlonsoMarcosM/TFM_Alonso_Marcos_Mu-oz}} Se asume Windows 10 u 11 con PowerShell, Docker Desktop, Node.js 20 o superior, Python 3.11 o superior y MiKTeX para compilar la memoria. Los comandos están extraídos de \texttt{scripts/infra/} y de la documentación canónica del repositorio.

\section{Preparación del entorno}

\begin{lstlisting}[language=PowerShell, caption={Instalación de dependencias Python del núcleo \texttt{om\_dcat\_sync}.}, label=code:ap-pip-install, captionpos=b]
python -m pip install --upgrade pip
python -m pip install -e .\tfm_ingestor
\end{lstlisting}

\begin{lstlisting}[language=PowerShell, caption={Instalación de dependencias de la consola web con pnpm.}, label=code:ap-pnpm-install, captionpos=b]
cd .\web
pnpm install
\end{lstlisting}

\section{Despliegue de infraestructura}

\begin{lstlisting}[language=PowerShell, caption={Lanzamiento de Kind, Helm e ingesta doble.}, label=code:ap-launch-infra, captionpos=b]
powershell -ExecutionPolicy Bypass -File .\scripts\infra\launch_infra.ps1
powershell -ExecutionPolicy Bypass -File .\scripts\infra\ingest_postgres_double.ps1
\end{lstlisting}

\section{Ejecución del workflow canónico}

\begin{lstlisting}[language=PowerShell, caption={Ejecución del workflow canónico end-to-end.}, label=code:ap-workflow, captionpos=b]
python -m om_dcat_sync workflow run --dry-run
python -m om_dcat_sync workflow run --allow-warnings
\end{lstlisting}

\section{Suite reproducible}

\begin{lstlisting}[language=PowerShell, caption={Suite reproducible documentada.}, label=code:ap-suite, captionpos=b]
powershell -ExecutionPolicy Bypass -File .\scripts\infra\run_validation_suite.ps1
\end{lstlisting}

La salida de la suite incluye \fqn{runtime_validation_report.json}, \fqn{dcat_catalog.jsonld} y \fqn{dcat_validation_report.ttl} en \fqn{tmp_pytest/}.
